\begin{appendix}

\section{Proofs and Source Verification}\label{app:proofs}

\begin{proof}[Proof of Theorem~\ref{thm:opportunity}]
Fix a representation \(R\) and component \(g\), and condition throughout on
\((\mathcal D_n,\mathcal Z_n)\). Put
\(\mathcal A_n^\Theta=\sigma(\mathcal D_n,\Theta,\mathcal Z_n)\). Uniformly in
the integrated norm, the conditional score decomposition gives
\begin{align*}
\E\{\Gamma_n(Y',W',X')\mid X'\in L_n^R(x),\mathcal A_n^\Theta\}
&=\mu+\sum_{h=1}^G
\E\{f_h(X')\mid X'\in L_n^R(x),\mathcal A_n^\Theta\}+o_p(1).
\end{align*}
Combining this identity with
\eqref{eq:honest-score-leaf} and \eqref{eq:component-leaf-reduction} gives the
component leaf reduction in \eqref{eq:semantic-leaf-reduction}.

Let \(X\sim P_X\) be independent of the training sample and draw a fresh
potential failure path conditional on \((X,\mathcal D_n,\mathcal Z_n)\). Write
\begin{equation*}
P_{g,n}^R(X)=
\prod_{t=1}^{J_{g,n}^{R,0}(X)}
\{1-\rho_{g,n,t}^{R,0}(X)\}.
\end{equation*}
Equation~\eqref{eq:survival-product} follows from the chain rule for the
sequential conditional failure kernels. Entry into the candidate set alone does
not determine success, and no threshold-localization approximation enters this
identity. If a reached history has hazard one, its factor is zero and the path
moves to the absorbing state, so the product still equals the zero conditional
failure probability.

Write \(\Lambda=\Lambda_{g,n}^R(X)\) and
\(\bar\rho=\bar\rho_{g,n}^R(X)\). On \(\{\bar\rho<1\}\),
\begin{align}
0
&\le -\log P_{g,n}^R(X)-\Lambda\notag\\
&=\sum_{t=1}^{J_{g,n}^{R,0}(X)}
\{-\log(1-\rho_{g,n,t}^{R,0}(X))-\rho_{g,n,t}^{R,0}(X)\}\notag\\
&\le
\frac{\bar\rho}{1-\bar\rho}\Lambda.
\label{eq:app-log-survival-bound}
\end{align}
On \(\{\lambda_g^R(X)<\infty\}\), conditions
\eqref{eq:hazard-limit-condition} and
\eqref{eq:app-log-survival-bound} imply
\(P_{g,n}^R(X)\to_p e^{-\lambda_g^R(X)}\). On
\(\{\lambda_g^R(X)=\infty\}\),
\[
0\le P_{g,n}^R(X)\le e^{-\Lambda_{g,n}^R(X)}\to_p0.
\]
Thus, under the joint law of \(X\) and the potential path,
\(P_{g,n}^R(X)-e^{-\lambda_g^R(X)}\to_p0\). The difference is bounded, so it
also converges to zero in mean. The failure-path identity and Jensen's
inequality therefore give
\[
\E\!\int\left|
\E_\Theta\{1-D_{g,n}^R(x)\mid\mathcal D_n,\mathcal Z_n\}
-e^{-\lambda_g^R(x)}
\right|dP_X(x)
\le \E|P_{g,n}^R(X)-e^{-\lambda_g^R(X)}|\longrightarrow0.
\]
Markov's inequality therefore gives
\begin{equation}\label{eq:app-conditional-survival}
\int\left|
\E_\Theta\{1-D_{g,n}^R(x)\mid\mathcal D_n,\mathcal Z_n\}
-e^{-\lambda_g^R(x)}
\right|dP_X(x)=o_p(1).
\end{equation}

Because \(g\) was arbitrary, \eqref{eq:app-conditional-survival}
holds for every \(h\in\{1,\ldots,G\}\).
Let \(\overline T_n^R(x)=
\E_\Theta\{\widetilde T_n^R(x)\mid\mathcal D_n,\mathcal Z_n\}\). By
the component leaf reduction, Jensen's inequality, and
\eqref{eq:app-conditional-survival},
\begin{equation}\label{eq:app-infinite-forest-target}
\left\|\overline T_n^R-
\mu-\sum_{h=1}^G\{1-e^{-\lambda_h^R(\cdot)}\}f_h
\right\|_{L^1(P_X)}=o_p(1).
\end{equation}
Conditional on \((\mathcal D_n,\mathcal Z_n)\), partition the trees into
mutually independent groups of size at most \(L<\infty\), with \(L=1\) in the
conditionally independent case. Conditional covariances vanish across groups.
Within a group, Cauchy--Schwarz and the bounded second moment
bound every covariance by a constant. Hence
\[
\E\!\left[
\left\|\frac1{B_n}\sum_{b=1}^{B_n}\widetilde T_{b,n}^R
-\overline T_n^R\right\|_{L^2(P_X)}^2
\middle|\mathcal D_n,\mathcal Z_n\right]\le\frac{CL}{B_n}=o(1).
\]
The \(L^1\)-norm is bounded by the \(L^2\)-norm. Combining this fact with
\eqref{eq:general-forest-linearization} and
\eqref{eq:app-infinite-forest-target} proves
\eqref{eq:general-opportunity-target}. The endpoint statements follow when the
limit function \(\lambda_g^R(\cdot)\) is identically zero, finite and positive,
or infinite.
\end{proof}

\begin{proof}[Proof of Theorem~\ref{thm:adaptive}]
All probability statements below are
uniform over \(a\in\mathcal C_a\). Constants may depend on
\((\underline a,\overline a,\Delta_a,\beta)\), but not on \(a\) or \(n\).
Write \(\mathcal D_n^+=\sigma(\mathcal D_n,\mathcal Z_n)\). Conditional
source-proof statements hold the nuisance-training folds, fitted nuisance
functions, and their training randomization fixed by conditioning on
\(\mathcal D_n^+\).

For every tree root, generate the subsample and honesty allocation first. The
ordinary tree uses independent uniforms indexed by node address to realize the
source transition law. For each representation \(h\in\{1,2\}\), construct a full
failure tree \(\mathcal T_{h,n}^0\) on the same persistent root samples. At a
reached node, a fresh node-address uniform realizes the regular conditional
source transition given that singleton \(X_{\bar h}\) is not selected. This
conditioning depends on the node state but not on an evaluation point. Both
children are then generated recursively. An independent \(X\sim P_X\) is drawn
only after the full tree is constructed and is routed through it. Every actual
or potential node is consequently an axis-aligned rectangle on a fixed root
sample. One relative-VC event can cover all of these rectangles without
conditioning on a selected target path. Indexing iid source draws by node
address is a distributional construction and makes no claim about the
package's sequential pseudorandom-number stream under a numerical seed.

Fix a node \(C\), condition on \(Q_M=q\), and rank the classes as in
Section~\ref{sec:resolution-intensity}. Class \(g_j(C)\) wins exactly
when the candidate set contains no column from classes
\(g_1(C),\ldots,g_{j-1}(C)\) and at least one
column from class \(g_j(C)\). Hence
\[
\Prb\{g_j(C)\text{ wins}\mid Q_M=q,C\}
=\frac{\binom{M-M_{j-1,C}^{\rm cum}}q
-\binom{M-M_{j,C}^{\rm cum}}q}{\binom Mq},
\]
because the two terms in the numerator count, respectively, the candidate
sets that exclude the preceding classes and those that also exclude class
\(g_j(C)\). This proves \eqref{eq:winner-fixed}. Conditioning on \(Q_M\)
proves \eqref{eq:winner-random}. The same count gives
\[
\Prb\{g\text{ is absent}\mid Q_M=q\}
=\frac{\binom{M-r_g}q}{\binom Mq},
\]
which proves \eqref{eq:availability-random}.

In either representation \(R_{h,n}\), the other coordinate \(X_{\bar h}\) is a
singleton. Hence
\begin{equation}\label{eq:app-singleton}
\Prb\{X_{\bar h}\text{ is included}\mid Q_{M_n}=q\}
=1-\frac{\binom{M_n-1}q}{\binom{M_n}q}
=\frac q{M_n}.
\end{equation}
Taking expectations gives the singleton's per-node inclusion probability
\(\omega_n=\E Q_{M_n}/M_n\).

To relate the source and population split criteria, recall that version 2.4.0
implements scalar \nolinkurl{causal_forest} fits through the
\nolinkurl{instrumental_trainer} path \citep{grfsource240}. The source class is
\nolinkurl{InstrumentalRelabelingStrategy}. Under \eqref{eq:production-dgp},
define
\[
\Gamma=4(W-1/2)\{Y-m_a(X)\},
\qquad
\widehat\Gamma=4(W-1/2)\{Y-\widehat m(X)\}.
\]
Random assignment and \(\E(U\mid X,W)=0\) give
\begin{equation}\label{eq:app-gamma-mean}
\E(\Gamma\mid X=x)=\tau_a(x),
\end{equation}
and uniform nuisance consistency gives
\begin{equation}\label{eq:app-score-nuisance}
\|\widehat\Gamma-\Gamma\|_\infty
\le2\|\widehat m-m_a\|_\infty=o_p(1).
\end{equation}
For a fresh observation, conditioning on the nuisance-training state therefore
yields
\begin{equation}\label{eq:app-conditional-score-mean}
\left\|\E(\widehat\Gamma\mid X=\cdot,\mathcal D_n^+)-\tau_a\right\|_\infty
=o_p(1).
\end{equation}

Set \(\ell_n=\log(B_ns_n)\) and
\[
\epsilon_n^{\rm split}
=\sqrt{\ell_n/k_n}+\ell_n/k_n+\|\widehat m-m_a\|_\infty,
\qquad
\epsilon_n^{\rm leaf}=\sqrt{\ell_n/k_n}.
\]
Both rates vanish under Assumption~\ref{ass:source-regularity}. Axis-aligned
rectangles, candidate children, and their intersections with treatment arms
form a fixed finite union of VC classes. Relative VC Bernstein bounds, followed
by a union bound over the \(B_n\) splitting and estimation samples, apply
simultaneously to every actual or potential node. Raw column multiplicity does
not enlarge these geometric classes. Since every admissible child contains at
least \(k_n\) observations from each treatment arm, the same event gives,
uniformly over all admissible children \(A\),
\[
\left|
\frac{\mathbb P_n^{\rm split}(\widehat\Gamma\one_A)}
{\mathbb P_n^{\rm split}(A)}-\E(\Gamma\mid A)
\right|=O_p(\epsilon_n^{\rm split}).
\]
Consequently, the CART criterion based on \(\widehat\Gamma\) converges uniformly
to its population counterpart \(\mathcal Q_C(j,t)\).

It remains to compare this criterion with the source criterion. Direct expansion
of \eqref{eq:source-relabel} gives
\[
4\Psi_{i,C}=\widehat\Gamma_i-\widehat\theta_C+E_{i,C},
\]
where
\[
E_{i,C}=-4Z_i\bar{\widetilde Y}_C-4\bar Z_C\widetilde Y_i
+4\bar Z_C\bar{\widetilde Y}_C
+8\widehat\theta_CZ_i\bar Z_C-4\widehat\theta_C\bar Z_C^2.
\]
The same VC event, applied to \(Z\), \(Y-m_a(X)\), and the nuisance error,
implies
\[
\sup_C\{|\bar Z_C|+|\bar{\widetilde Y}_C|\}
=O_p(\epsilon_n^{\rm split}),
\qquad
\sup_C|\widehat\theta_C|=O_p(1).
\]
Thus \(\sup_C\max_{i\in C}|E_{i,C}|=O_p(\epsilon_n^{\rm split})\). The term
\(-\widehat\theta_C\) is constant within the parent and cancels from every
child-mean contrast. Expanding the two bounded split criteria now gives
\begin{equation}\label{eq:app-source-population}
\sup_{C,j,t}
|\widehat{\mathcal Q}_C^{\rm grf}(j,t)-\mathcal Q_C(j,t)|
=O_p(\epsilon_n^{\rm split})=o_p(1).
\end{equation}

Consider coordinate \(j\) before its first selection on a path. Its projection in the
current rectangular node is \([0,1]\). Conditional independence makes all
other effect components cancel from a split contrast. For a split at \(t\),
the population criterion is
\begin{equation}\label{eq:app-signed-cart-shape}
\mathcal Q_C(j,t)=
\begin{cases}
\displaystyle a_j^2\frac{t}{1-t},&0<t\le1/2,\\[6pt]
\displaystyle a_j^2\frac{1-t}{t},&1/2\le t<1.
\end{cases}
\end{equation}
It has unique maximizer \(1/2\) and maximum \(a_j^2\). The gain gap between
the two unresolved coordinates is at least
\(\Delta_a(2\underline a+\Delta_a)>0\).
Set
\[
\delta_*=\frac18\min\{\underline a^2,
\Delta_a(2\underline a+\Delta_a)\}.
\]
By \eqref{eq:app-source-population}, the uniform criterion error is below
\(\delta_*\) with probability tending to one. This event preserves both the
unresolved-coordinate ordering and the gap between an untouched signal and the
residual gain of a previously selected signal.

Fix \(\eta>0\). At a node with \(n_C\ge(4+\eta)k_n\), the relative VC bounds
make a sample threshold \(1/2+o_p(1)\) source admissible uniformly. If
\(4k_n\le n_C<(4+\eta)k_n\), every admissible child has at least \(2k_n\)
observations. Both empirical child fractions therefore lie in
\([2/(4+\eta),1-2/(4+\eta)]\), and uniform convergence of the untouched
coordinate's conditional empirical distribution forces every admissible
threshold to be within \(O(\eta)+o_p(1)\) of \(1/2\). The positive argmax gap
in \eqref{eq:app-signed-cart-shape}, \eqref{eq:app-source-population}, and the
order of limits \(n\to\infty\) followed by \(\eta\downarrow0\) give
\begin{equation}\label{eq:app-first-split-localization}
\max_{\mathcal T,j,F}
|\widehat t_{j,F}-1/2|=o_p(1),
\end{equation}
where \(j\in\{1,2\}\), \(\widehat t_{j,F}\) is the threshold selected when
coordinate \(j\) is first used at node \(F\), and \(\mathcal T\) ranges over
every actual tree and potential failure tree. The maximum is zero when the set
is empty.

To control later splits after one coordinate has been selected, let
\(b_j(C)=\min\{\Prb(\iota_j=0\mid C),\Prb(\iota_j=1\mid C)\}\) for a
rectangular node \(C\). Conditional independence and total variance imply
\[
\sup_t\mathcal Q_C(j,t)\le4a_j^2b_j(C),
\qquad
\sum_{\ell=1}^KP_X(C_\ell)b_j(C_\ell)
\le P_X(C)b_j(C)
\]
for every rectangular partition \(C_1,\ldots,C_K\) of \(C\). The second
inequality follows by summing the smaller of the unconditional masses for
\(\iota_j=0\) and \(\iota_j=1\) in each child.

Let \(\mathcal N_{j,\mathcal T}\) be the disjoint nodes at which coordinate
\(j\) is first selected on branches of tree \(\mathcal T\), and set
\begin{equation}\label{eq:app-weighted-threshold-error}
\mathcal E_{j,\mathcal T}
=\sum_{F\in\mathcal N_{j,\mathcal T}}
P_X(F)|\widehat t_{j,F}-1/2|.
\end{equation}
Since \(X_j\mid F\sim\operatorname{Unif}[0,1]\), the two children created at
\(F\) have total probability-weighted \(b_j\)-impurity equal to
\(P_X(F)|\widehat t_{j,F}-1/2|\). Thus every descendant frontier
\(\mathcal P\) satisfies
\begin{equation}\label{eq:app-frontier-residual-gain}
\sum_{C\in\mathcal P}P_X(C)\sup_t\mathcal Q_C(j,t)
\le4\overline a^2\mathcal E_{j,\mathcal T}.
\end{equation}
The first-split nodes form an antichain. Hence
\(\mathcal E_{j,\mathcal T}\le\max_F|\widehat t_{j,F}-1/2|\le1/2\).
Equation \eqref{eq:app-first-split-localization}, bounded convergence, and
iterated expectation give
\begin{equation}\label{eq:app-weighted-error-vanishes}
\E_\Theta(\mathcal E_{j,\mathcal T}\mid\mathcal D_n^+)=o_p(1),
\qquad j=1,2.
\end{equation}
The same conclusion holds for either potential failure tree because
\eqref{eq:app-first-split-localization} is uniform over those trees.

Set \(\gamma=\underline a^2/4\). After coordinate \(j\) has been selected,
let \(\mathcal N_{j,\mathcal T}^{\rm hit}(\gamma)\) contain the first descendant
node on each branch where another coordinate remains admissible and
\(\sup_t\mathcal Q_C(j,t)>\gamma\). These nodes form an antichain, so
\eqref{eq:app-frontier-residual-gain} gives
\begin{equation}\label{eq:app-first-hitting-bound}
\sum_{C\in\mathcal N_{j,\mathcal T}^{\rm hit}(\gamma)}P_X(C)
\le\frac{4\overline a^2}{\gamma}\mathcal E_{j,\mathcal T}.
\end{equation}
An untouched signal has an admissible near-half split with gain at least
\(\underline a^2+o(1)\). Outside these descendant subtrees, a previously
selected signal has gain at most \(\gamma\). The uniform bound
\eqref{eq:app-source-population} then makes the untouched signal the preferred
admissible signal with probability tending to one.

Let \(\mathcal U_n\) collect the uniform criterion, admissibility, and
localization events above, so \(\Prb(\mathcal U_n^c)=o(1)\). For representation
\(h\), let \(\mathcal B_{h,n}^0(x)\) contain \(\mathcal U_n^c\) and the event
that the potential failure path of \(x\) enters one of the first-hitting
subtrees. Equations \eqref{eq:app-weighted-error-vanishes} and
\eqref{eq:app-first-hitting-bound} imply
\begin{equation}\label{eq:app-adaptive-resolution}
\max_{h\in\{1,2\}}
\E_{\Theta_h^0}\!\left[
\Prb_X\{\mathcal B_{h,n}^0(X)\mid\Theta_h^0,\mathcal D_n^+\}
\middle|\mathcal D_n^+\right]=o_p(1).
\end{equation}
Outside \(\mathcal B_{h,n}^0(X)\), every admissible unselected signal included
in the candidate set defeats each signal selected earlier on the path, except
at the initial competing-component split and the fixed-width terminal boundary.

On paths where signal \(j\) is first selected, let
\(\widehat\iota_{j,\mathcal T}\) be its empirical side label and let
\(\mathcal R_{j,\mathcal T}\) be the union of terminal cells below such splits.
The disagreement mass at node \(F\) equals
\(P_X(F)|\widehat t_{j,F}-1/2|\), and therefore
\begin{equation*}
\E_\Theta\Prb_X\{X\in\mathcal R_{j,\mathcal T},
\widehat\iota_{j,\mathcal T}(X)\ne \iota_j(X)\}=o_p(1).
\end{equation*}
For an independent \(X'\), write \(L_{\mathcal T}(x)\) for the terminal cell
containing \(x\). Inserting the leaf-constant empirical label and applying
conditional Jensen gives
\begin{align}
&\int_{\mathcal R_{j,\mathcal T}}
\left|
\E\{\iota_j(X')\mid X'\in L_{\mathcal T}(x),\mathcal T\}-\iota_j(x)
\right|dP_X(x)\notag\\
&\qquad\le
2\Prb_X\{X\in\mathcal R_{j,\mathcal T},
\widehat\iota_{j,\mathcal T}(X)\ne \iota_j(X)\}.
\label{eq:app-leaf-contamination}
\end{align}
Thus imperfect threshold localization contributes \(o_p(1)\) in integrated
loss.

We now control the number of split opportunities. Fix \(h\) and route \(x\)
through the full failure tree
\(\mathcal T_{h,n}^0\). Its transitions use the original law over \(M_n\)
columns conditional on failure to select \(X_{\bar h}\). No column is deleted,
and \(Q_{M_n}\) is never redrawn under a different dimension. The increasing
recodings of \(X_h\) induce identical routing at corresponding thresholds. Let
\(n_d(x)\) be the number of
splitting observations in the node containing \(x\) after \(d\) splits.

We establish the depth envelope on the uniform relative-VC event. The
information restriction and uniform treatment balance imply that every
admissible child obeys, for some fixed \(\beta_0\in(0,\beta)\),
\[
\beta_0n_d(x)-1\le n_{d+1}(x)
\le(1-\beta_0)n_d(x)+1
\]
simultaneously over all actual and failure-tree nodes. To justify continuation,
consider any rectangular node with \(n_d(x)\ge C_{\rm cont}k_n\), where
\(C_{\rm cont}\) is a sufficiently large fixed constant. Its treatment share is
\(1/2+o(1)\),
uniformly over roots and rectangles. Continuity of the covariate distribution
makes the observations in each coordinate distinct almost surely, so every
sampled coordinate is nonconstant whenever the node contains at least two
observations. Its empirical median creates
two children of size \(n_d(x)/2+O(1)\), and each child contains
\(n_d(x)/4+o\{n_d(x)\}\) observations from each treatment arm. For a sufficiently
large fixed \(C_{\rm cont}\), these counts exceed \(k_n\). Moreover,
\[
\frac{n_A(1/4-\bar Z_A^2)}
{n_d(x)(1/4-\bar Z_C^2)}=\frac12+o(1)>\beta
\]
for both median children. The median of every sampled column is therefore
source admissible. Since \(Q_{M_n}\ge1\), the candidate set contains such a
split. Condition~\eqref{eq:source-split-nondegeneracy} then forces a positive
criterion value and continuation. Nodes with
\(n_d(x)\ge C_{\rm cont}k_n\) therefore cannot be
terminal. No split is admissible when \(n_d(x)<4k_n\), because each child must
contain at least \(k_n\) observations on both sides of the parent treatment
mean. The uncertain range
\(4k_n\le n_d(x)<C_{\rm cont}k_n\) has fixed multiplicative
width and contributes
only a bounded number of additional splits under the displayed recursion.
These statements hold on one event over every actual and full failure-tree node
because the root samples are fixed before the trees are generated. Choosing
\[
0<c_\beta<|\log\beta_0|^{-1},
\qquad C_\beta>|\log(1-\beta_0)|^{-1}
\]
uses the fastest permitted contraction \(\beta_0\) for the lower depth bound and
the slowest contraction \(1-\beta_0\) for the upper bound. These choices give,
simultaneously for \(h=1,2\),
\begin{equation}\label{eq:app-random-depth-envelope}
\Prb\{c_\beta H_n\le J_{h,n}^0(X)\le C_\beta H_n\}\to1.
\end{equation}
The same lower and upper
recursions apply to every actual path, regardless of which coordinate is
selected.

For fixed constants \(0<c_{\rm leaf}<C_{\rm leaf}<\infty\), every reached
terminal cell \(L\) has population mass
\begin{equation}\label{eq:app-leaf-mass}
c_{\rm leaf}\frac{k_n}{s_n}\le P_X(L)\le
C_{\rm leaf}\frac{k_n}{s_n}
\end{equation}
with probability tending to one. For the upper bound, a node above a
sufficiently large multiple of \(k_n\) admits a median split, and
\eqref{eq:source-split-nondegeneracy} forces continuation. The lower bound
follows from the treatment-arm and information restrictions on the split that
created the leaf. Consequently there are at most
\(c_{\rm leaf}^{-1}s_n/k_n\) leaves per tree.

For the estimation half, apply the relative VC inequalities to every
axis-aligned rectangle. Write
\[
N_{b,n}^{\rm est}(L)
=\sum_{i\in I_{b,n}^{\rm est}}\one\{X_i\in L\}.
\]
For \(L=L_{b,n}^R(x)\), this count equals
\(N_{b,n}^{R,{\rm est}}(x)\) from Section~\ref{sec:framework}.
Together with \eqref{eq:app-leaf-mass}, these inequalities imply, for fixed
constants \(0<c_{\rm est}<C_{\rm est}<\infty\),
\begin{equation*}
c_{\rm est}k_n\le N_{b,n}^{\rm est}(L)\le C_{\rm est}k_n
\end{equation*}
for every nonempty honest leaf, with probability tending to one. Honesty
pruning is therefore asymptotically inactive. Write
\(\widehat\Gamma_i=\Gamma_i+\delta_{i,n}\), where
\(\max_i|\delta_{i,n}|\le2\|\widehat m-m_a\|_\infty\) by
\eqref{eq:app-score-nuisance}. Conditional on the splitting tree, the oracle
scores \(\Gamma_i\) in its independent estimation half are independent across
observations and satisfy \(\E(\Gamma_i\mid X_i)=\tau_a(X_i)\). Applying the
relative-VC bound to these oracle scores and
then adding the uniform nuisance remainder gives
\begin{equation}\label{eq:app-honest-moment}
\max_{b,L}\left|
\frac1{N_{b,n}^{\rm est}(L)}
\sum_{i\in L\cap I_{b,n}^{\rm est}}\widehat\Gamma_i
-\E\{\tau_a(X)\mid X\in L\}
\right|
=O_p\!\left(\sqrt{\frac{\ell_n}{k_n}}
+\|\widehat m-m_a\|_\infty\right)=o_p(1).
\end{equation}
For such a leaf, write
\[
\bar Z_{b,n}(L)
=\frac{1}{N_{b,n}^{\rm est}(L)}
\sum_{i\in L\cap I_{b,n}^{\rm est}}Z_i.
\]
The same argument for \(Z\) gives
\begin{equation}\label{eq:app-honest-treatment-balance}
\max_{b,L}|\bar Z_{b,n}(L)|=O_p(\epsilon_n^{\rm leaf})=o_p(1).
\end{equation}
The forest-weighted \(\bar Z_\alpha(x)\) is an average of these leaf means.
Boundedness of \(\widetilde Y\), \eqref{eq:source-prediction-remainder}, and
\eqref{eq:app-honest-treatment-balance} therefore yield
\begin{equation}\label{eq:app-source-prediction-transfer}
\max_{h\in\{1,2\}}
\|\widehat T_{h,n}-\check T_{h,n}\|_\infty
=O_p(\epsilon_n^{\rm leaf})=o_p(1).
\end{equation}

It remains to derive the target for each representation. Work under the full
failure-tree law used to define
\(\rho_{h,n,t}^0\). Condition on the training sample and the history through the
current potential state, and write \(\bar h=3-h\).
The singleton is available with probability \(\omega_n\) by
\eqref{eq:app-singleton}, independently of the past. It may enter the candidate
set without being chosen for the split. Let \(K_{h,n}^{\rm str}(x)\) count internal split nodes
at which an available singleton can fail because a higher-gain unresolved
component wins or the node lies in the final fixed-width leaf-size boundary.

For \(h=1\), the singleton has the larger unresolved gain and therefore wins
whenever it is available, outside the exceptional set. For \(h=2\), the repeated
component has the larger unresolved gain. A split at which that component
blocks the singleton also resolves the repeated component. Once that split has
been taken, an
available unresolved singleton defeats every previously selected component on
\(\mathcal U_n\) outside the exceptional subtrees identified in the
adaptive-resolution argument. The failure-tree depth recursion leaves at most
a fixed number \(C_{\rm bd}\) of internal
nodes in the terminal boundary. Hence
\(K_{h,n}^{\rm str}(X)\le 1+C_{\rm bd}\) on
\(\mathcal U_n\) for both representations. Since
\(\mathcal U_n^c\subseteq\mathcal B_{h,n}^0(X)\) has probability \(o(1)\), the
random-depth envelope gives
\begin{equation}\label{eq:app-structural-node-share}
\max_{h\in\{1,2\}}
\frac{K_{h,n}^{\rm str}(X)}{J_{h,n}^0(X)}\longrightarrow_p0
\end{equation}
in integrated probability, with the ratio defined as zero when
\(J_{h,n}^0(X)=0\).

For \(t\le J_{h,n}^0(X)\), the predictable difference
\(\omega_n-\rho_{h,n,t}^0(X)\) is nonnegative. At a structural split node, this
difference is at most \(\omega_n\). At every other internal split node
on the complement of \(\mathcal B_{h,n}^0(X)\), the adaptive-resolution
argument shows statewise that singleton availability is equivalent to
successful singleton resolution, so
the difference is zero. Therefore, under the same potential failure-transition law,
\begin{equation*}
0\le J_{h,n}^0(X)\omega_n-\Lambda_{h,n}^0(X)
\le K_{h,n}^{\rm str}(X)\omega_n
+J_{h,n}^0(X)\omega_n\one\{\mathcal B_{h,n}^0(X)\}.
\end{equation*}
Every term in this pathwise bound is evaluated under the same potential
failure-transition law, so no probability from the ordinary path law is
substituted. Dividing by the
normalization \(1+J_{h,n}^0\omega_n\) gives, with the ratio interpreted
as zero when \(J_{h,n}^0=0\),
\begin{align}
&\frac{J_{h,n}^0(X)\omega_n-\Lambda_{h,n}^0(X)}
{1+J_{h,n}^0(X)\omega_n}\notag\\
&\qquad\le
\frac{K_{h,n}^{\rm str}(X)}{J_{h,n}^0(X)}
\frac{J_{h,n}^0(X)\omega_n}{1+J_{h,n}^0(X)\omega_n}
+\one\{\mathcal B_{h,n}^0(X)\}
\le
\frac{K_{h,n}^{\rm str}(X)}{J_{h,n}^0(X)}
+\one\{\mathcal B_{h,n}^0(X)\}.\notag
\end{align}
The first term is \(o_p(1)\) by
\eqref{eq:app-structural-node-share}, and the second is \(o_p(1)\) under the
joint law of \(X\) and the potential failure-transition randomizers by
\eqref{eq:app-adaptive-resolution}. Consequently,
\begin{equation}\label{eq:app-source-hazard-equivalence}
\max_{h\in\{1,2\}}
\frac{\left|\Lambda_{h,n}^0(X)-J_{h,n}^0(X)\omega_n\right|}
{1+J_{h,n}^0(X)\omega_n}\longrightarrow_p0,
\qquad
\max_{h\in\{1,2\}}\max_{1\le t\le J_{h,n}^0(X)}
\rho_{h,n,t}^0(X)\le\omega_n,
\end{equation}
where the inner maximum is zero when \(J_{h,n}^0(X)=0\).
Iterated conditioning on the potential failure states gives the exact
failure-path identity in \eqref{eq:survival-product}.

The depth and failure-path statements in this proof use the capped-Poisson
source law \eqref{eq:grf-q} with finite \(\xi_{M_n}\). A candidate law that
samples all \(M_n\) columns almost surely requires a different probability
calculation.
It remains to verify, including when \(M_n\) is fixed, that the repeated
component is resolved before the actual path reaches its leaf. Its conditional
candidate-inclusion probability is
\[
p_n^{\rm rep}=1-\frac{\nu_{M_n}(1)}{M_n}\ge\frac12.
\]
The actual-path versions of the antichain and depth arguments show, with
conditional probability \(1-o_p(1)\), that the path avoids the exceptional
subtrees and has at least \(\lfloor c_\beta H_n\rfloor\) internal nodes. Set
\(L_n=\lfloor c_\beta H_n\rfloor-C_{\rm bd}\). On this event, the first
inclusion of the repeated class either resolves it or loses to and resolves the
higher-gain singleton. Its next inclusion resolves the repeated component.
Successive conditioning on the fresh node-address draws therefore bounds the
failure probability by
\[
\Prb\{\operatorname{Bin}(L_n,p_n^{\rm rep})\le1\}
\le(1+2L_n)2^{-L_n}=o(1).
\]
Let \(\mathcal B_{h,n}^{\rm cop}(X)\) denote failure to resolve the repeated
component before the leaf. The preceding bounds give
\begin{equation}\label{eq:app-copied-before-leaf}
\max_{h\in\{1,2\}}
\E_\Theta\!\left[
\Prb_X\{\mathcal B_{h,n}^{\rm cop}(X)\mid\Theta,\mathcal D_n^+\}
\middle|\mathcal D_n^+\right]=o_p(1),
\end{equation}
including fixed \(M_n\).

For a generic tree \(\mathcal T\), write
\[
m_{j,\mathcal T}(x)
=\E\{s_j(X')\mid X'\in L_{\mathcal T}(x),\mathcal T\}.
\]
If the singleton is unresolved, the terminal cell is a function of \(X_h\)
outside the integrated exceptional event. Independence and centering then give
\(m_{\bar h,\mathcal T}(x)=0\). If a component is resolved,
\eqref{eq:app-leaf-contamination} controls the difference between its leaf mean
and its sign at \(x\). Together with \eqref{eq:app-copied-before-leaf}, these
facts give the componentwise bounds
\begin{align}
\E_\Theta\!\left[
\|m_{h,\mathcal T}-s_h\|_{L^1(P_X)}\mid\mathcal D_n^+\right]
&=o_p(1),\notag\\
\E_\Theta\!\left[
\|m_{\bar h,\mathcal T}
-D_{\bar h,n}^{R_{h,n}}s_{\bar h}\|_{L^1(P_X)}
\mid\mathcal D_n^+\right]&=o_p(1).
\label{eq:app-componentwise-leaf}
\end{align}
The losses are bounded, so the joint fresh-tree bounds imply the displayed
conditional statements by the tower property and Markov's inequality. Combining
\eqref{eq:app-componentwise-leaf} with \eqref{eq:app-honest-moment} gives
\begin{equation}\label{eq:app-tree-mixture}
\E_\Theta\!\left[
\left\|
\widetilde T_n^{R_{h,n}}
-D_{\bar h,n}^{R_{h,n}}\tau_a
-\{1-D_{\bar h,n}^{R_{h,n}}\}a_hs_h
\right\|_{L^1(P_X)}\middle|\mathcal D_n^+\right]=o_p(1).
\end{equation}
Equations \eqref{eq:app-gamma-mean}--\eqref{eq:app-conditional-score-mean}
verify Assumption~\ref{ass:semantic-score} with
\(\mu=0\), \(f_1=a_1s_1\), and \(f_2=a_2s_2\). Equations
\eqref{eq:app-honest-moment}, \eqref{eq:app-componentwise-leaf}, and
\eqref{eq:app-tree-mixture} verify the honest-score and component-reduction
parts of Assumption~\ref{ass:honest-resolution}. The bounded independent tree
groups in Assumption~\ref{ass:source-regularity}, together with
\eqref{eq:app-source-prediction-transfer}, verify its reported-forest
linearization.

The repeated component is already handled directly by
\eqref{eq:app-copied-before-leaf}. No separate failure kernel is needed for it.
For the singleton, the survival-product argument in the proof of
Theorem~\ref{thm:opportunity} applies to \(\Lambda_{h,n}^0\). Combining its
conditional resolution probability with the mixture
\eqref{eq:app-tree-mixture}, forest aggregation, and
\eqref{eq:app-source-prediction-transfer} gives the source-forest target.

If \(\mathcal I_n\to0\), the upper depth envelope
\eqref{eq:app-random-depth-envelope} and
\eqref{eq:app-source-hazard-equivalence} give
\(\Lambda_{h,n}^0(X)\to_p0\), so the singleton resolution probability vanishes
and \eqref{eq:signed-selection-targets} follows. If
\(\mathcal I_n\to\infty\), the lower envelope gives
\(\Lambda_{h,n}^0(X)\to_p\infty\), so the singleton is resolved with probability
tending to one. Together with \eqref{eq:app-copied-before-leaf}, this proves
\eqref{eq:signed-recovery-target}. For a fixed \(a\in\mathcal C_a\) in the
finite-intensity regime,
\(\omega_n\to0\), and \eqref{eq:intermediate-depth} together with
\eqref{eq:app-source-hazard-equivalence} gives
\[
\Lambda_{h,n}^0(X)\xrightarrow[\Prb_a^0]{p}\kappa d_h(a),
\qquad
\max_{1\le t\le J_{h,n}^0(X)}\rho_{h,n,t}^0(X)
\xrightarrow[\Prb_a^0]{p}0.
\]
The survival-product argument yields singleton failure probability
\(\zeta_h(a)=e^{-\kappa d_h(a)}\). Substitution into
\eqref{eq:app-tree-mixture} gives the limits
\eqref{eq:intermediate-target-one}--\eqref{eq:intermediate-target-two}.

The same coupling covers the increasing recodings used in the theorem. At a
node, the source represents a threshold split by a column value \(t\) and sends
an observation left when its value is at most \(t\). For every
\(\varphi_{g\ell n}\), the corresponding transformed threshold is
\(\varphi_{g\ell n}(t)\), and
\(\varphi_{g\ell n}(x)\le\varphi_{g\ell n}(t)\) exactly when \(x\le t\).
The candidate partitions, admissibility checks, split gains, and routing are
therefore identical within each recoding block. The proof above depends on a
block only through its multiplicity, so all three limits remain unchanged.
\end{proof}

\begin{proof}[Proof of Corollary~\ref{cor:sign-policy}]
In the selection regime, for \(g\in\{1,2\}\),
Markov's inequality and the margin \(a_g\ge\underline a\) give
\begin{equation}\label{eq:app-sign-classification}
\Prb_X\{\operatorname{sgn}(\widehat T_{g,n}(X))\ne s_g(X)\}
\le\frac{2}{\underline a}
\|\widehat T_{g,n}-a_gs_g\|_{L^1(P_X)}=o_p(1).
\end{equation}
Since \(s_1,s_2\) are independent Rademacher variables,
\(\Prb(s_1s_2<0)=1/2\). On this event,
\(|\tau_a(X)|=a_2-a_1\ge\Delta_a\). The triangle inequality and another use of
Markov's inequality show that, except on a set of \(P_X\)-mass \(o_p(1)\),
\[
\min\bigl\{|\widehat T_{1,n}(X)|,|\widehat T_{2,n}(X)|\bigr\}
\ge\underline a/2.
\]
This proves
\eqref{eq:population-sign-reversal} and
\eqref{eq:strong-population-sign-reversal}.

Because \(a_2>a_1>0\), \(\operatorname{sgn}\{\tau_a(X)\}=s_2(X)\). Equation
\eqref{eq:app-sign-classification} therefore proves
\eqref{eq:oracle-sign-policy}. For the other representation,
\[
\operatorname{Reg}_a(\one\{s_1=1\})
=\E\left[\tau_a
\{\one\{s_2=1\}-\one\{s_1=1\}\}\right]
=\frac{a_2-a_1}{2}.
\]
The bounded outcome and policy-classification convergence transfer this value
to \(\widehat\pi_{1,n}\), proving \eqref{eq:sign-policy-regret}.
\end{proof}

\begin{proof}[Proof of Corollary~\ref{cor:critical-multiplicity}]
For every \(M\), the capped-Poisson rule satisfies the exact identity
\[
\E Q_M=\xi_M+e^{-\xi_M}-\E(N_M-M)_+.
\]
If \(M\to\infty\) and \(\xi_M\equiv\xi\), the last expectation converges to
zero, which proves \eqref{eq:fixed-mtry-boundary}. If \(\xi_M\to\infty\) and
\(\xi_M/M\to0\), Poisson concentration makes the upper cap irrelevant and
\(\E Q_M/\xi_M\to1\). Under \eqref{eq:default-mtry},
\(\xi_M/\sqrt M\to1\), which proves \eqref{eq:default-phase}. A fixed
finite representation has \(\omega_n>0\) constant and
\(H_n\omega_n\to\infty\), so it is in the recovery regime.
\end{proof}

\begin{proof}[Proof of Proposition~\ref{prop:finite-g}]
Because \(G\) is fixed, axis-aligned rectangles and coordinate halfspaces in
\([0,1]^G\) have fixed VC dimension. The criterion-transfer, localization,
first-hitting, honest-leaf, and source-prediction arguments in the proof of
Theorem~\ref{thm:adaptive} therefore hold uniformly over
\(a\in\mathcal C_G\), all coordinates, and all \(G\) representations. The split
shape remains \eqref{eq:app-signed-cart-shape}, and
\eqref{eq:finite-g-coefficients} gives a uniform ordering gap among unresolved
coordinates. Outside an integrated event of probability \(o(1)\), an included
unresolved coordinate with the largest gain is selected, its first threshold
is \(1/2+o_p(1)\), and its leaf contamination is \(o_p(1)\).

Fix \(R_{g,n}^{(G)}\). Each singleton has conditional inclusion probability
\(\omega_{G,n}=\E Q_{M_{G,n}}/M_{G,n}\), while the repeated class has probability
\[
p_{g,n}^{\rm rep}
=1-\E\!\left\{
\frac{\binom{G-1}{Q_{M_{G,n}}}}
{\binom{M_{G,n}}{Q_{M_{G,n}}}}
\right\}
\ge \frac{c_n}{c_n+G-1}\ge\frac1G.
\]
In the selection regime, the path has at most \(C_\beta H_n\) internal nodes
with probability tending to one. The probability that some singleton enters a
candidate set on the path is at most
\((G-1)C_\beta H_n\omega_{G,n}+o(1)=o(1)\).
On its complement, all splits use the repeated class. Independence, centering,
and the uniform leaf and prediction transfers give
\(\|\widehat T_{g,n}^{(G)}-a_gs_g\|_{L^1(P_X)}=o_p(1)\). Maximizing over fixed
\(G\) proves \eqref{eq:finite-g-selection}.

For recovery, the path has at least \(c_\beta H_n\) internal nodes with
probability tending to one. Divide them into \(G\) consecutive blocks of length
\(L_n\asymp H_n/G\). Fresh node-address draws and the preceding bound imply,
for every split-action class \(j\), that its probability of no inclusion in a
given block is at most \(\exp(-L_n\omega_{G,n})=o(1)\).
Order coordinates by decreasing \(a_j^2\). Successive blocks resolve the
largest-gain remaining coordinate outside the common exceptional event. A union
bound over fixed \(G\) shows that all coordinates are resolved before stopping,
which proves \eqref{eq:finite-g-recovery}.

For \(g\ne h\), \eqref{eq:finite-g-selection} and the margin
\(a_j\ge\underline a_G\) imply
\(\Prb_X\{\operatorname{sgn}(\widehat T_{j,n}^{(G)})\ne s_j\}=o_p(1)\) for
\(j\in\{g,h\}\).
Independence gives \(\Prb(s_gs_h<0)=1/2\), proving
\eqref{eq:finite-g-pair-sign}. The multiplicity boundaries follow from the
calculation for Corollary~\ref{cor:critical-multiplicity} after replacing
\(M_n\) by \(M_{G,n}=c_n+G-1\).
\end{proof}

\begin{proof}[Proof of Proposition~\ref{prop:quotient}]
Couple two quotient-preserving representations by all the common randomization
listed before the proposition. At any common node state, they draw the same
classes, evaluate the same certified actions, apply the same scores and tie
order, and select the same action. The weak-order certificate gives identical
unordered children, and canonical orientation routes every point in
\(\mathcal X_{\rm cert}\) to the same child. Induction gives identical tree
partitions on the certified array. Leaf observations and forest weights then
coincide. The measurability condition extends this equality to nuisance fits,
stopping decisions, sufficient statistics, predictions, reported variance
calculations, rankings, and policy ties. This proves the pathwise statement.

At node \(C\), class \(g\) wins when it is sampled and none of the
\(h_g(C)\) classes preceding it in the total score order is sampled. Conditional
on \(\widetilde Q_K=q\),
\[
\Prb\{g\text{ wins}\mid \widetilde Q_K=q,C\}
=\frac{\binom{K-h_g(C)}q-\binom{K-h_g(C)-1}q}{\binom Kq}.
\]
The two terms in the numerator count candidate sets that exclude the preceding
classes, with and without class \(g\). Averaging over \(q\) proves
\eqref{eq:quotient-hazard}.
\end{proof}

\end{appendix}
